% ---------------------------------------------------------------------------
% DRAFT for the reviewer pass. Not generated by paper/build_all.py.
%
% Replaces / extends the second paragraph of Section 5. Suggested section
% rename: "Discussion, Limitations and Conclusion" -- the content is already
% there, but reviewers scan for the word.
%
% Each claim below is backed by a number in output/run_final/analysis/.
% ---------------------------------------------------------------------------

\paragraph{Limitations.}
Four limitations bound what our results establish.

First, the node axis changes the target as well as the graph. A secondary-entity
node has no category of its own, so its label is the modal category over the
primary records it participates in. Differences along the node axis therefore
combine a construction effect with a change of prediction problem, and
cross-node-type comparisons should be read with that in mind. We report the
node population, class count and class prior for every pair in
Table~\ref{tab:nodestats}; notably, the priors do not explain the effects we
measure. ArXiv's two node types have near-identical priors ($20.9\%$ and
$24.3\%$) yet differ by more than twenty accuracy points, and Electronics N2
has the more favourable prior of the two ($51.5\%$ against $36.5\%$) yet scores
roughly thirty points lower. The node-type effect is therefore not an artifact
of class imbalance.

Second, our proxy subsets are not scale models of the full graph. Edge rules
that match on an attribute produce a node degree that grows with pool size,
while $k$-nearest-neighbour rules do not: going from $1{,}000$ rows to the full
History pool multiplies edges per node by $17.8\times$ for E4b and $17.9\times$
for E4c but leaves E5b unchanged at $1.01\times$. The same asymmetry makes
primary-entity graphs very sparse at subset scale --- $90.8\%$ of ArXiv N1 nodes
are isolated in a $1{,}000$-row subset against $8.4\%$ of N2 nodes --- so part
of the advantage we measure for secondary entities on ArXiv, and for
similarity-based edges generally, reflects connectivity supplied by the
sampling rather than by the construction. Section~\ref{sec:fullscale} quantifies
how much of the subset ranking survives at full scale.

Third, the decision tree is an interpretability layer rather than a predictor.
Ranking variants by their raw train-pool mean achieves the same held-out
correlation as the fitted tree ($\rho = 0.939$ against $0.934$ for node
classification; $0.858$ against $0.875$ for retrieval), so the rank-correlation
results in Table~\ref{tab:dt_consistency} establish that proxy scores are
reproducible across pools, not that the meta-learner adds predictive power. The
tree's own contribution is that it compresses the score table into a small
number of readable rules and that it extends to constructions it never scored,
which we measure by leave-one-variant-out ($\rho = 0.802$ and $0.692$).

Fourth, every score comes from one architecture, one sentence encoder, and one
parameterization per edge rule. The E5 arms in particular are single points in
a continuum: replacing $k=50$ with $k=5$ in E5b left GNN accuracy statistically
unchanged but moved the RAGAS composite by $0.28$, which exceeds the
best-to-worst variant spread we report on several datasets. We therefore regard
the specific edge rankings as conditional on these parameterizations, while the
finding that construction choice matters, and that its effect is
dataset-dependent, does not rest on them.
